\section{总体分析}

本文先回答数据如何评价，再比较资源投入的效果，最后讨论预算分配与能力变化。四个问题使用的数据、方法及相互联系见图\ref{fig:overall_architecture}。

\begin{figure}[htbp]
  \captionsetup{font={bundledsong,minusfour,bf}}
  \centering
  % 总体架构：横向箭头对应每问的数据、方法和产出。
% 不在此图中表达尚未核实的跨问数值校准关系。
\begingroup
\definecolor{archblue}{HTML}{315A79}
\definecolor{archlight}{HTML}{EDF3F7}
\definecolor{archgray}{HTML}{536572}
\begin{tikzpicture}[
  font=\song\fontsize{10}{13}\selectfont,
  archbox/.style={draw=archblue!55, line width=0.55pt, rounded corners=2pt,
    fill=archlight, align=center, inner sep=7pt, minimum height=1.55cm},
  archdata/.style={archbox, text width=2.55cm},
  archmodel/.style={archbox, text width=5.1cm},
  archoutput/.style={archbox, text width=3.3cm, fill=white},
  archlink/.style={-{Stealth[length=2mm]}, draw=archblue, line width=0.7pt},
  archhead/.style={text=archblue, font=\song\bfseries\fontsize{11}{14}\selectfont}
]
\node[archhead] at (-5.8,0.65) {数据与条件};
\node[archhead] at (-0.8,0.65) {建模任务与方法};
\node[archhead] at (4.8,0.65) {主要产出};

\node[archdata] (a) at (-5.8,-0.75) {附件 A\\质量信号\\配方与验证损失};
\node[archmodel, fill=archblue, text=white] (q1) at (-0.8,-0.75)
  {问题一：数据质量与领域配比\\熵权评价、冲突分析、岭回归};
\node[archoutput] (o1) at (4.8,-0.75)
  {综合质量评分\\配比与损失关系\\推荐配比};

\node[archdata] (b) at (-5.8,-3.0) {附件 B\\规模、训练量\\与质量信息};
\node[archmodel] (q2) at (-0.8,-3.0)
  {问题二：损失与资源的关系\\广义标度律、边际效用与弹性};
\node[archoutput] (o2) at (4.8,-3.0)
  {损失函数与参数\\资源之间的替代关系};

\node[archdata] (c) at (-5.8,-5.25)
  {算力预算\\题给成本设定\\C7 上下文信息};
\node[archmodel] (q3) at (-0.8,-5.25)
  {问题三：算力约束下的配置\\成本建模与约束优化};
\node[archoutput] (o3) at (4.8,-5.25)
  {最优配置 $(N,D,Q)$\\预算变化下的\\配置特征};

\node[archdata] (d) at (-5.8,-7.5)
  {附件 C\\历史模型\\与评测信息};
\node[archmodel] (q4) at (-0.8,-7.5)
  {问题四：能力演进与前沿预测\\损失与评测桥接、回归与外推};
\node[archoutput] (o4) at (4.8,-7.5)
  {规模与非规模贡献\\能力前沿预测区间};

\foreach \archsrc/\archmid/\archdst in {a/q1/o1,b/q2/o2,c/q3/o3,d/q4/o4} {
  \draw[archlink] (\archsrc.east) -- (\archmid.west);
  \draw[archlink] (\archmid.east) -- (\archdst.west);
}
\end{tikzpicture}
\endgroup

  \caption{算力约束下资源配置建模的总体架构}
  \label{fig:overall_architecture}
\end{figure}

问题一从语料开始。我们用熵权法合成质量指标，通过成对标准化得分差识别分歧，用 Kendall 协同系数检查整体一致性，再以岭回归描述领域占比与损失的关系。由此得到质量评分 $Q$、基线 $Q_0$ 和配比回归系数，供后续分析使用。

问题二比较不同投入能使损失降低多少。在 Chinchilla 形式上加入质量缺口，得到 $L(N,D,Q)=E+AN^{-\alpha}+BD^{-\beta}+C(1-Q)^{\gamma}$；当 $Q=1$ 时，质量项归零，恢复经典函数形式。随后对参数、数据和质量求导，并比较等损失条件下的替代关系。

问题三在上述损失函数中加入预算限制。成本分别计算为基础训练 $6ND$、质量提升 $D[g(Q)-g(Q_0)]_+$ 和注意力开销 $\eta NDL_{\text{ctx}}$，在给定配比情景下用 SLSQP 求解。除比较不同预算的配置外，还以 $L_{\text{crit}}=6/\eta$ 说明两类训练成本何时相等，并从预算扫描中识别质量投入的阶段变化。

问题四转向下游评测。先按数据可比程度拟合损失与得分，再用含年份固定效应的面板回归分解历史增量，以 logistic 曲线和结构情景预测前沿。C8 的子任务结果用于寻找能力短板，C4 的宏观记录用于检查资源变化及数据覆盖。

四问的衔接主要体现在变量和结果的传递上：问题一的 $Q_0$ 与配比修正 $M(\bm p)$、问题二估计的标度律参数，共同用于问题三。问题四检查损失与评测得分是否同向变化；由于两者的映射误差较大，后续预测直接使用评测得分，而不将前三问的损失收益精确换算为能力分数。这样既保留了各问之间的联系，也区分了训练损失与实际评测所能说明的内容。
